% answers_short.tex —— 2011 年·安徽卷·文科 · 简短答案速查
% 由 tools/make_short_answers.py 自 answers.tex 的 \gksolhead/\gkanswer 提取，
% 并与 2011/2011_数学真题_安徽文_答案要点.md 逐题比对后生成，勿手改（改源 MD 后重跑）。

\begin{gknotice}
\textbf{使用说明}：本节只给最终答案便于快速核对，完整推导见解析版。答案为编写者独立推导，\textbf{非官方标准答案}。
\end{gknotice}

\gkpart{选择题}
\begin{center}\begin{tabular}{ccccccccccc}
\toprule
题号 & 1 & 2 & 3 & 4 & 5 & 6 & 7 & 8 & 9 & 10 \\
\midrule
答案 & A & B & C & B & D & B & A & C & D & A \\
\bottomrule
\end{tabular}\end{center}

\gkpart{填空题}
\begin{itemize}[leftmargin=2.6em]
  \item[\textbf{11.}] $-3$。
  \item[\textbf{12.}] $15$。
  \item[\textbf{13.}] $(-3,2)$。
  \item[\textbf{14.}] $\dfrac{\pi}{3}$。
  \item[\textbf{15.}] ①③。
\end{itemize}

\gkpart{解答题}
\begin{itemize}[leftmargin=2.6em]
  \item[\textbf{16.}] $\dfrac{\sqrt3+1}{2}$。
  \item[\textbf{17.}] (1) 证明见解析；(2) 证明见解析
  \item[\textbf{18.}] (1) 极大值点 $x=\dfrac12$，极小值点 $x=\dfrac32$；(2) $0<a\leqslant1$
  \item[\textbf{19.}] (1) 证明见解析；(2) $\dfrac32$
  \item[\textbf{20.}] (1) $\hat y=6.5x-12778.8$（$x$ 为年份）；(2) $299.2$ 万吨
  \item[\textbf{21.}] (1) $a_n=n+2$；(2) $S_n=\dfrac{\tan(n+3)-\tan3}{\tan1}-n$
\end{itemize}
